% Problem record op_e8307b66d76af92a. This TeX file is derived from
% database/problems_json/op_e8307b66d76af92a.json by scripts/sync-tex.mjs; edit the JSON record
% and rerun the script rather than editing this file.
\section{One-way state generation from private Hamiltonian phase states}
\label{sec:op_e8307b66d76af92a}

\paragraph{Problem.}
Do private-architecture Hamiltonian phase states yield a one-way state generator for explicit polynomial parameters?

A function $\varepsilon(n)$ is negligible if, for every constant $a>0$, eventually $\varepsilon(n)<n^{-a}$. “QPT” denotes a quantum polynomial-time algorithm.

Question. Are there explicit polynomially bounded parameter functions $m(n)$ and growing $q(n)$ for which the private-architecture Hamiltonian-phase-state ensemble is a one-way state generator?

For a uniformly sampled $A\in\mathbb F_2^{m\times n}$ and independently uniform $\theta_i\in\{2\pi j/q:0\leq j<q\}$, define

\begin{equation}
|\phi_{A,\theta}\rangle
=\exp\!\left(i\sum_{i=1}^{m}\theta_i Z^{A_i}\right)|+\rangle^{\otimes n},
\qquad Z^{A_i}:=\bigotimes_{j=1}^{n}Z^{A_{ij}}.
\label{eq:e830-1}
\end{equation}

This is the Hamiltonian-phase-state architecture studied in \sourcecite{ref:e830-bostanci25}{Bostanci25} and subsequently in \sourcecite{ref:e830-cao26}{Cao26}. The concrete claim is that, for every polynomial $t$ and every QPT algorithm $\mathcal I$, with

\begin{equation}
(A',\theta')\leftarrow
\mathcal I\!\left(1^n,|\phi_{A,\theta}\rangle^{\otimes t(n)}\right),
\label{eq:e830-2}
\end{equation}

one has

\begin{equation}
\mathbb E\!
\left[\left|\langle\phi_{A',\theta'}|\phi_{A,\theta}\rangle\right|^2\right]
\leq\operatorname{negl}(n).
\label{eq:e830-3}
\end{equation}

The output must be a valid description in the same public parameter space; invalid outputs fail verification. The architecture $A$ is hidden. This is a parameter-explicit existence formulation of search HPS, not a claim that every parameter choice is secure. \sourcecite{ref:e830-bostanci25}{Bostanci25}, Sections~4.1

The displayed definitions, constraints, and target bounds are recorded in Eqs.~\eqref{eq:e830-1}, \eqref{eq:e830-2}, \eqref{eq:e830-3}.

\subsection*{Status}

Unsolved

\subsection*{Source}

This precise formulation is editor wording based on the unresolved direction and limitations documented in the cited primary literature \sourcecite{ref:e830-bostanci25}{Bostanci25}; it is not presented as a verbatim conjecture of those authors.

\subsection*{Progress}

\begin{itemize}
  \item Status: an unresolved quantum cryptographic hardness assumption. The originating work gives restricted worst-to-average-case reductions and bounded-copy design evidence, but not a proof of general search hardness. It also explains that polynomially many copies suffice information-theoretically: the conjectured obstacle is computation, not an absence of information. \sourcecite{ref:e830-bostanci25}{Bostanci25}, Sections~5.1–5.3

  \item Later work develops measurement-assisted shallow preparation of these states and analyzes statistical properties. These are advances in realizing the ensemble, not proofs that polynomial-time quantum inversion is impossible. \sourcecite{ref:e830-cao26}{Cao26}

  \item The cited work leaves both general inversion and security unresolved. The original paper separately proposes decision HPS, concerning indistinguishability from Haar states; that different security task is not identified here with search HPS. \sourcecite{ref:e830-bostanci25}{Bostanci25}, Sections~4.2
\end{itemize}

\subsection*{References}

\begin{enumerate}[label={},leftmargin=4.8em,labelsep=0.5em]
  \footnotesize
  \item[\textup{[Bostanci25]}]\label{ref:e830-bostanci25}
  John Bostanci, Jonas Haferkamp, Dominik Hangleiter, and Alexander Poremba, \emph{Efficient Quantum Pseudorandomness from Hamiltonian Phase States}. \href{https://arxiv.org/html/2410.08073}{arXiv:2410.08073}; \href{https://drops.dagstuhl.de/entities/document/10.4230/LIPIcs.TQC.2025.9}{TQC 2025, LIPIcs 350, article 9}. Pinpoints refer to the arXiv full text: §4.1, §4.2, and §§5.1–5.3.

  \item[\textup{[Cao26]}]\label{ref:e830-cao26}
  Chenfeng Cao and Jens Eisert, \emph{Measurement-Driven Quantum Advantages in Shallow Circuits}. \href{https://arxiv.org/html/2505.04705}{arXiv:2505.04705}; Physical Review Letters 136, 080601 (2026). See “Measurement-prepared Hamiltonian phase states” and the discussion.
\end{enumerate}

\subsection*{Comment}

This is the closest match to a quantum-information version of a one-way function: a short classical description prepares a state, but copies allegedly do not permit efficient reverse engineering.

The fidelity verifier matters. Recovering a different description of almost the same state is a successful attack; merely showing that the original labels are nonunique is not evidence of security. Conversely, tomography with an exponentially expensive reconstruction stage does not refute a polynomial-time hardness claim.

Scope caution: the security parameter regime is part of the research problem. The statement does not endorse an arbitrary choice such as $m=n$, nor a version in which the architecture is public. The cited results establish partial evidence rather than security for an explicit general parameter regime.

\subsection*{Field}

Quantum Cryptography

\subsection*{Topic}

Computational complexity and computability; Quantum state preparation

\subsection*{ID}

\texttt{op\_e8307b66d76af92a}
